\documentclass{article}

\usepackage{booktabs}
\usepackage{tabularx}
\usepackage{microtype}

\title{Development Plan\\\progname}

\author{\authname}

\date{}

\input{../Comments.text}
\input{../Common.text}

\begin{document}

\maketitle

\begin{table}[hp]
\caption{Revision History} \label{TblRevisionHistory}
\begin{tabularx}{\textwidth}{llX}
\toprule
\textbf{Date} & \textbf{Developer(s)} & \textbf{Change}\\
\midrule
Sep 26th 2026 & Omar + Meena & Section 1-5\\
Sep 27th 2026 & Meena & Section 6/7/8\\
Sep 27th 2026 & Marko & Section 9/11/13\\
Sep 27th 2026 & Omar & Section 10/12/Appendices\\
\bottomrule
\end{tabularx}
\end{table}

\newpage{}

\wss{Put your introductory blurb here.  Often the blurb is a brief roadmap of
what is contained in the report.}

\wss{Additional information on the development plan can be found in the
\href{https://gitlab.cas.mcmaster.ca/courses/capstone/-/blob/main/Lectures/L02b_POCAndDevPlan/POCAndDevPlan.pdf?ref_type=heads}
{lecture slides}.}

\section{Confidential Information?}

There is no confidential information from industry to protect.
\section{IP to Protect}

There is no external IP partnership to protect. This project is independently thought of and implemented by the student team.

\section{Copyright License}

Our team is not adopting a copyright license as of now, therefore our team's work falls under default copyright laws.

\section{Team Meeting Plan}

The team will have a weekly in-person meeting; those with extenuating circumstances may join virtually. The default meeting time is Thursday at 12:30pm as none of the team members have lectures. \newline
\newline
We plan to meet with our supervisor every ~2 weeks on a Zoom call, subject to their availability. Larger development milestones will ideally be showcased in-person to the supervisor. \newline
\newline
The meetings will be structured in a manner where the Project Manager comes prepared with an agenda of upcoming deliverables and tasks to be handled. The Project Manager is the meeting chair facilitating the flow of the meeting, such as major updates from each team member. Outstanding work or conflicts are addressed before the end of the meeting. The agenda denotes all decisions made regarding agenda items.



\section{Team Communication Plan}

The team will use Git Issues with specific labels on a Kanban Board display to track the progress of work items and who is responsible for them. Github branches and pull requests function as a form of version control, history and technical reviews. \newline
\newline
Teams will be utilized as the formal method of communication for file sharing, important messages and virtual meetings. An iMessage groupchat serves as an informal means to deliver quick non-sensitive information.

\section{Team Member Roles}

\wss{You should identify the types of roles you anticipate, like notetaker,
leader, meeting chair, reviewer.  Assigning specific people to those roles is
not necessary at this stage.  In a student team the role of the individuals will
likely change throughout the year.}

\section{Workflow Plan}

\begin{itemize}
	\item How will you be using git, including branches, pull request, etc.?
	\item How will you be managing issues, including template issues, issue
	classification, etc.?
  \item Use of CI/CD
\end{itemize}

\section{Project Decomposition and Scheduling}

\begin{itemize}
  \item How will you be using GitHub projects?
  \item Include a link to your GitHub project
\end{itemize}

\wss{How will the project be scheduled?  This is the big picture schedule, not
details. You will need to reproduce information that is in the course outline
for deadlines.}

\section{Proof of Concept Demonstration Plan}

\noindent Risks:
\begin{itemize}
  \item Different users may find some perspectives not conducive to learning or engaging; for example, being overwhelmed by the nature of 3D perspectives.

  \noindent What we will demonstrate to convince ourselves we can overcome this risk is designing multiple perspective modes for suitable games (e.g., 2D, first-person 3D, third-person 3D) to alleviate any user issues pertaining to perspective.

  \noindent If our PoC demo fails regarding this, we may work around this risk by supplanting the 3D ideas with 2.5D rendering instead, replicating games such as Overcooked.

  \item Users may find some games to be overly stimulating, or not stimulating enough; for example, the noise level in a grocery shopping mini-game.

  \noindent What we will demonstrate to convince ourselves we can overcome this risk is adding adjustable settings to each game to pinpoint the most suitable level of stimulation for the user.

  \noindent If our PoC demo fails regarding this, we may simplify the simulation levels to a binary option, or assume low stimulation as the default, since overstimulation is a larger risk.

  \item Some games may be considered too easy or too difficult depending on the user, missing the goal of learning and progress tracking.

  \noindent What we will demonstrate to convince ourselves we can overcome this risk is implementing difficulty levels for each game. This enables users to play at a difficulty ideal for their skills and needs.

  \noindent If our PoC demo fails regarding this, we may implement a binary option regarding difficulties that solely shortens the maximum time given to complete a task within the game, or the game as a whole.

  \item The caretaker may not have sufficient flexibility or freedom to tailor games to each user; the games not being personable per user.
  
  \noindent What we will demonstrate to convince ourselves we can overcome this risk is implementing a template system for the caretaker to create their own games from pre-existing baseline games. For example, the caretaker may create a vocabulary-based mini-game for the user's specific information such as their favorite grocery items.

  \noindent If our PoC demo fails regarding this, we may create a text-based form per game for the caretaker to input the customized vocabulary, rather than a full template system with pictures, colors et cetera.

  \item The caretaker may have difficulty interpreting the data collected from the games, and thus be unable to derive useful information for the user.
  
  \noindent What we will demonstrate to convince ourselves we can overcome this risk is creating multiple views of the data for the caretaker such as graphical views, table views and summarized bullet point views. 

  \noindent If our PoC demo fails regarding this, we may instead offer a method for the caretaker to export the data to a CSV file, and analyze it themselves in a program such as Excel.
\end{itemize}

\par\medskip

\section{Expected Technology}

Our team will use a pre-packaged game engine, namely Godot, as the bedrock of our project implementation. Godot is a free and open-source game engine that has a very lenient learning curve via a Pythonic language (GD Script) made specifically for game development. It is also cross-platform, allowing for enhanced accessibility as the project can run on multiple devices.

\par\medskip

\noindent AssetLib, Godot's native asset library, enables us to choose assets and add-ons such as: objects \& frameworks instead of creating everything from scratch. Godot Unit Testing (GUT) is a unit testing framework that serves as a quick method for creating and executing unit tests.

\par\medskip

\noindent For frontend purposes, we will largely be using Godot's built-in GUI system, which allows for the creation of user interfaces with minimal code. However, we may use the TauPlot add-on for complex data visualization instead of creating our own math system.

\par\medskip

\noindent For backend purposes, we are looking at Godotbase as a bridge between our UI layer and potentially our Supabase database layer. Godotbase permits very easy integration with Supabase, while Supabase has strong features for authentication, database management and storage. 

\par\medskip

\noindent SonarQube is a code quality tool that we may use to double as both a linter and a code coverage measuring tool. It is compatible with Godot and can be used to analyze our codebase for potential bugs, vulnerabilities, and code smells.

\par\medskip

\noindent Our plan for Continuous Integration (CI) is to use GitHub Actions, which is a CI/CD platform that allows us to automate our build, test, and deployment processes. It is integrated with GitHub, making it easy to set up and manage.


\section{Use of AI and LLMs}

\begin{itemize}

\item \textbf{Development and Documentation:} LLMs will be used throughout development as a helpful resource for both programming and documentation. They will assist with becoming familiar with specific and game-engine functionality by providing explanations, examples, and implementation guidance.
\newline \textbf{Verification:} The team will focus on testing generated code strictly.
\item \textbf{Code Generation:} LLMs may be used to generate boilerplate, unit tests, and common code structures to speed up development and avoid wasting time, allowing us to focus on design.
\newline \textbf{Verification:} All generated code will be reviewed by the team to fit with the project's requirements, and go through unit testing and playtesting.
\item \textbf{Content Bank Generation:} Using AI to generate large banks of game-specific content from a predefined ruleset (such as vocabulary exercises, puzzles, situational problems, etc.) can greatly help with a game's educational value, by encouraging players to learn the skill rather than memorize sets and answers. 
\newline \textbf{Verification:} Outputs will be manually reviewed and tested to ensure that they make sense, and align with the specific learning objectives of the given game.
\item \textbf{Asset Generation:} For audio/visual development, non LLM based tools like stable diffusion and 3D mesh-generation may be used to produce specific sprites, UI assets, or mesh assets if needed for devopment speed. 
\newline \textbf{Verification:} The generated assets will have to follow a given visual style guide, for clarity and accessibility.

\end{itemize}

\section{Coding Standard}

The official GDScript style guide will be followed for all code written in GDScript. This includes things such as naming conventions, code element ordering, and other best practices. The team will also follow industry standard coding habits such as writing clear and concise comments, using meaningful variable and function names, and adhering to the principles of clean code.

\par\medskip

\noindent The official style guide can be found at \url{https://docs.godotengine.org/en/stable/getting_started/scripting/gdscript/gdscript_styleguide.html}.


\newpage{}

\section{Appendix --- Generative AI Use Disclosure}



\subsection{AI Tools Used}

Claude's Sonnet 5 model.

\subsection{How and Where Used}

Sonnet was used to learn some of the quirks of LaTeX and fix errors associated with it. This includes the bullet point structure in the Team Member Roles section and paragraphs in the PoC Demo Plan.


\subsection{Verification of AI-Generated Content}

AI-Generated content was not used in this document. The usage was limited to learning about LaTeX and fixing associated errors.

\newpage{}

\section*{Appendix --- Reflection}

\input{../Reflection.text}

\begin{enumerate}
    \item Why is it important to create a development plan prior to starting the
    project?

    \item In your opinion, what are the advantages and disadvantages of using
    CI/CD?
    \item What disagreements did your group have in this deliverable, if any,
    and how did you resolve them?
\end{enumerate}

\newpage{}

\section*{Appendix --- Team Charter}

\subsection*{External Goals}

Our team has an aligned set of goals for this project. We want to create a project that has a real positive impact on people's lives, especially a marginalized group such as autistic individuals, because our aim is to use our acquired skills to leave the world a better place than we found it. 

\par \medskip

\noindent Secondary priorities include winning a department podium award (1st, 2nd, or 3rd place) at the Capstone EXPO, and showcasing a complex project that we can be proud of to potential employers.

\subsection*{Attendance}

\subsubsection*{Expectations}

\noindent Punctuality and attendance is expected and required for all team members. If a team member is unable to attend a meeting, they must notify the team at least 30 minutes in advance. If a team member is late to a meeting, they must notify the team as soon as possible. If a team member needs to leave a meeting early, they must notify the team prior to the meeting if possible. If a team member is unable to attend multiple meetings in a row, they must provid a valid reason for their absence.


\subsubsection*{Acceptable Excuse}

For deadlines, only critical emergencies can function as an excuse for missing them. This is because the team member should notify other team members as soon as they think they might not be able to meet a deadline, so the team has time to find a solution and distribute the outstanding work.

\par \medskip

\noindent For one-time meeting absences, a team member reserves the right to keep their reason for absence private. However, a repeated absence requires a reason for absence and must constitute a reasonable excuse by the rest of the team that deterred them from attending.

\subsubsection*{In Case of Emergency}

Team members are expected to notify the team as soon as possible if they have an emergency so the team can redistribute their work accordingly. If appropriate, the team member should brief the team on the current status of their work, allowing others to pick up where they left off more smoothly.

\subsection*{Accountability and Teamwork}

\subsubsection*{Quality} 

\noindent Team members should be prepared to discuss their work and progress at each meeting. The meeting chair's agenda along with the Kanban board will reflect the current status of deliverables and work items; each team member is expected to provide updates regarding issued assigned to them. 

\par\medskip

\noindent The quality of deliverables should be strong with a focus on efficiency and clean code. Technical debt is natural and expected, but must be explicitly documented and addressed in a timely manner. Shortcuts and workarounds are only acceptable if they do not degrade the quality of the codebase, or if they are temporary solutions that will be addressed in the near future.
\subsubsection*{Attitude}

\noindent We expect all team members to exhibit openness to others' ideas no matter what their own ideas, in the name of respect and collaboration. Cooperation and positive attitudes are an expection, and any team member who acts otherwise will be addressed by the team to ensure everyone feels comfortable. 

\par\medskip

\noindent For conflicts, it is expected to be resolved between the concerned parties. Polite communication is imperative. If a conflict still remains unresolved, the team members may utilize one or more other team members as mediators to aid in the conflict resolution. The team may then elect to take a vote or bring the issue to the TA or instructor.
\subsubsection*{Stay on Track}

The team has a meeting chair and project manager that together manage the team's progress; they will have the ability to tell if the team is off track. For team members, it is expected that they will complete their work in a timely manner and at a high level. For team members who are struggling to complete their work, the environment should be supportive enough to allow them to ask for help and guidance from other team members. Team members not contributing their fair share will be addressed by the team in a kind manner and encouraged to contribute proportionately if they are able to. If the problem persists, the team may elect to bring the issue to the TA or instructor.

\par \medskip

\noindent Meeting attendance and amount of lines of code committed to the repository will be monitorted by team members as a metric for contribution. The consequences of not doing proportionate work relative to others will be addressed internally, with empathy, before it becomes a persistent issue. However, a pattern of not contributing will be brought to the attention of the TA and/or instructor. Team members performing particularly well may be rewarded with priase and retain the right to choose what tasks they prefer to do in the following sprint or deliverable.


\subsubsection*{Team Building}

Our team is already well-acquainted with each other, so we already spend fun time together outside of the project and have a good rapport with respect and trust. We will continue to build team cohesion with activities such as eating meals together, playing games, and celebrating milestones. 

\subsubsection*{Decision Making} 

Decisions are made after each group member has spoken their opinion. It is important that each member feel valued and genuinely heard, not figuratively. Options will be discussed until a consensus is reached. Ideally, decisions are unanimous, but quick decisions may have to be taken with a majority vote.

\end{document}